\documentclass[12pt,a4paper,titlepage]{article}
\usepackage[latin1]{inputenc}
\usepackage{amsmath}
\usepackage{amsfonts}
\usepackage{amssymb}
\usepackage{graphicx}
\usepackage{listings}
\usepackage{datetime}
\usepackage{multirow}
\usepackage{hyperref}
\hypersetup{
    colorlinks,
    citecolor=black,
    filecolor=black,
    linkcolor=black,
    urlcolor=black
}
\newcommand{\HRule}{\rule{\linewidth}{0.5mm}}

\begin {document}

\begin{titlepage} 
\begin{center}

\includegraphics[width=1\textwidth]{./logos/freeems}\\[1cm] \vfill

\HRule \\[0.8cm]
{ \huge \bfseries Memory Management Scheme}\\[0.4cm]
\HRule \vfill

\Large 
%\emph{Author:}\\
%\textsc{Fred Cooke} \vfill

{\large Version 0.2 --- 27th November 2008}

\end{center}
\end{titlepage} 


\tableofcontents
\thispagestyle{empty}
\pagebreak


\section {Purpose / Preamble} 

This document exists to summarise and explain the memory available on the
\texttt{XDP512} processor and how it is being used in the FreeEMS application
firmware.

This document is believed to be true and correct at the time of writing.
Obviously the firmware is under continuous active development and there will
be changes to the way things are done as time passes. The code is the best
reference, however I will endeavour to keep this up to date.


\section{Overall Memory Layout}

The \texttt{XDP512} has a 16 bit data bus and main addressing scheme. This
limits addressable space to a total of 64k (65536 bytes). The 64k is made up of
the following parts:

\begin{itemize}
\item 2k of control registers from \texttt{0x0000-0x07FF} with 
length \texttt{0x0800}

\item 2k of addressable EEPROM from \texttt{0x0800-0x0FFF} with
length \texttt{0x0800}

\item 12k of addressable RAM from \texttt{0x1000-0x3FFF} with
length \texttt{0x3000}

\item 48k of addressable Flash from \texttt{0x4000-0xFFFF} with
length \texttt{0xC000}
\end{itemize}

The registers are permanently exposed and linearly addressed, and therefore
accessed directly at any time. The EEPROM, RAM and Flash, however, all use a
paging scheme to expose a larger space than would otherwise be possible. The
total amount of each accessible using paging is as follows:

\begin{itemize}
\item 4k of EEPROM

\item 32k of RAM

\item 512k of Flash
\end{itemize}


\section{EEPROM Memory Layout}

EEPROM is not used in FreeEMS at this time, however the 2k is made up of the
following parts:

\begin{itemize}
\item 1k of EEPROM page window from \texttt{0x0C00-0x0FFF} with
length \texttt{0x0400}

\item 1k of direct addressable EEPROM from \texttt{0x0800-0x0BFF}
with length \texttt{0x0400}
\end{itemize}

The following are valid \texttt{EPAGE} values: \texttt{0xFF}, \texttt{0xFE},
\texttt{0xFD}, \texttt{0xFC}. The fixed 1K page from \texttt{0x0C00} to
\texttt{0x0FFF} is the page exposed by setting \texttt{EPAGE-0xFF}.
Doing that is pointless however as you then have two copies of the same EEPROM
exposed at the same time. The reset value of \texttt{EPAGE} is \texttt{0xFE}
and ensures that there is a linear EEPROM space available between addresses
\texttt{0x0800} and \texttt{0x0FFF} out of reset. Therefore, the useful set of
\texttt{EPAGE} values are: \texttt{0xFE}, \texttt{0xFD}, \texttt{0xFC}


\section{RAM Memory Layout}

This section details how the available RAM is laid out and used in FreeEMS.
The 12k of direct addressable RAM consists of the following parts:
\begin{itemize}
\item 4k of RAM page window from \texttt{0x1000-0x1FFF} of length
\texttt{0x1000}

\item 8k of direct addressable RAM from \texttt{0x2000-0x3FFF} of
length \texttt{0x2000}
\end{itemize}

The following are valid \texttt{RPAGE} values: \texttt{0xFF, 0xFE, 0xFD, 0xFC,
0xFB, 0xFA, 0xF9, 0xF8}. The fixed 4K page from \texttt{0x2000-0x2FFF} is the
page exposed by setting \texttt{RPAGE} to \texttt{0xFE}. The fixed 4K page
from \texttt{0x3000-0x3FFF} is the page exposed by setting \texttt{RPAGE} to
\texttt{0xFF}. As with the EEPROM, doing that is pointless because you are just
exposing the region twice at the same time. The reset value of \texttt{RPAGE}
is \texttt{0xFD} and ensures that there is a linear RAM space available between
addresses \texttt{0x1000} and \texttt{0x3FFF} out of reset. Therefore the
useful set of \texttt{RPAGE} values is: \texttt{0xFD, 0xFC, 0xFB, 0xFA, 0xF9,
0xF8}.

In FreeEMS we use the 6 available pages as two triplets for tune switching
(like table switching, but almost everything is switched). At any given time
the firmware is using 3 pages that contain one copy of the following:

\begin{itemize}
\item 3 VE Tables

\item 1 Lambda Target Table

\item 2 Ignition Timing Tables

\item 2 Injection Timing Tables

\item 4 blocks of tunable items including all small tables etc.
\end{itemize}

Each of these 12 items is 1k in size totalling 12k which doubled (for tune
switching) is 24k. These are available in 3 sets of four blocks at a time.
VE and Lambda are together, Timing tables are together and the tunable items
are together. The remaining 8k is directly addressable linear space, used as
follows:

\begin{itemize}
\item 2k buffer for reception of communications data

\item 2k buffer for transmission of communications data

\item 4k of general purpose RAM
\end{itemize}

The 4k of general purpose RAM is used in the traditional way with all of the
global variables at the low end and the stack at the high end growing down
toward the globals. Currently there are approximately 1k of global variables in
that region.



\section{Flash Memory Layout}

This section details how the available Flash is laid out and used in FreeEMS.
The 48k of direct addressable Flash consists of the following parts:

\begin{itemize}
\item 16k of direct addressable Flash from \texttt{0x4000-0x7FFF}
of length \texttt{0x4000}

\item 16k of Flash page window from \texttt{0x8000-0xBFFF} of
length \texttt{0x4000}

\item 16k of direct addressable Flash from \texttt{0xC000-0xFFFF}
of length \texttt{0x4000}
\end{itemize}

The following are valid \texttt{PPAGE} values: \texttt{0xFF, 0xFE, 0xFD, 0xFC,
0xFB, 0xFA, 0xF9, 0xF8, 0xF7, 0xF6, 0xF5, 0xF4, 0xF3, 0xF2, 0xF1, 0xF0, 0xEF,
0xEE, 0xED, 0xEC, 0xEB, 0xEA, 0xE9, 0xE8, 0xE7, 0xE6, 0xE5, 0xE4, 0xE3, 0xE2,
0xE1, 0xE0}.

The fixed 16K page from \texttt{0x4000-0x7FFF} is the page exposed by setting
\texttt{PPAGE} to \texttt{0xFD}. The fixed 16K page from \texttt{0xC000-0xFFFF}
is the page exposed by setting \texttt{PPAGE} to \texttt{0xFF}. As with the RAM
and EEPROM, doing that is pointless because you are just exposing the region
twice at the same time.

The reset value of \texttt{PPAGE} is \texttt{0xFE} and ensures that there
is a linear Flash space available between addresses \texttt{0x8000} and
\texttt{0xBFFF} out of reset. Therefore, the useful set of \texttt{PPAGE}
values are: \texttt{0xFE, 0xFC, 0xFB, 0xFA, 0xF9, 0xF8, 0xF7, 0xF6, 0xF5, 0xF4,
0xF3, 0xF2, 0xF1, 0xF0, 0xEF, 0xEE, 0xED, 0xEC, 0xEB, 0xEA, 0xE9, 0xE8, 0xE7,
0xE6, 0xE5, 0xE4, 0xE3, 0xE2, 0xE1, 0xE0}.

The 512k of available flash space is made up of four 128k blocks, each with
eight 16k pages. Currently, due to limitations of hcs12mem and/or the serial
monitor, we can only use up to eight 16k pages totalling 128k of Flash. Two of
those pages are permanently exposed in linear space (\texttt{0xFF/text} and
\texttt{0xFD/text1}) leaving six pages of Flash for dynamic use in the short to
medium term.

The two linear regions are called \texttt{text} and \texttt{text1}, the six Flash pages are
swapped in and out automatically by the compiler whenever needed. In them we
store the following:

\begin{itemize}
\item Commonly used items and the serial monitor

\item All Interrupt Service Routines

\item Lookups and functions operating on them

\item Fuel tables and functions operating on them

\item Timing tables and functions operating on them

\item Tunable items and functions operating on them

\item General function storage area one

\item General function storage area two
\end{itemize}



\section{Where Things Are}

The current break down of usage within those pages is as follows:

Commonly used items and the serial monitor:
\begin{itemize}
\item All functions without a page specified are stored in text region by default
\item The serial monitor takes up 2k of space at the high end of text
\end{itemize}

All Interrupt Service Routines:
\begin{itemize}
\item All placed in text1 region as a way of tracking flash usage
\end{itemize}

Lookups and functions:
\begin{itemize}
\item 2k CHT lookup table
\item 2k IAT lookup table
\item 2k MAF lookup table
\item Core variable generation function
\end{itemize}

Fuel tables and functions:
\begin{itemize}
\item 6 1k VE tables
\item 2 1k Lambda tables
\item Copy flash to ram function
\end{itemize}

Timing tables and functions:
\begin{itemize}
\item 4 1k Ignition timing tables
\item 4 1k Injection timing tables
\end{itemize}

Tunable items and functions:
\begin{itemize}
\item 8 1k blocks of tunable items
\end{itemize}

General paged function space one:
\begin{itemize}
\item Assorted functions and data used purely within them
\end{itemize}

General paged function space two:
\begin{itemize}
\item Assorted functions and data used purely within them
\end{itemize}

Note: The ISR code \emph{must} be kept in linear space such that they can be
called asynchronously at any time.


\end {document}
